\documentclass[11pt,a4paper]{article}
\usepackage[utf8]{inputenc}
\usepackage[T1]{fontenc}
\usepackage[spanish]{babel}
\usepackage{geometry}
\geometry{a4paper, margin=0.8in}
\usepackage{xcolor}
\usepackage{parskip}

\definecolor{yariblue}{RGB}{29,58,107}
\definecolor{yarigreen}{RGB}{5,120,90}
\definecolor{yarigray}{RGB}{90,90,90}

\newcommand{\bloque}[3]{%
  \vspace{4pt}%
  \noindent{\color{yariblue}\bfseries #1}\hfill{\color{yarigreen}\small\itshape #2}\par%
  \vspace{1pt}%
  \noindent #3\par%
  \vspace{2pt}%
  {\color{yarigray}\hrulefill}\par%
}

\pagestyle{empty}

\begin{document}

\begin{center}
{\LARGE\bfseries\color{yariblue} YariNET --- Pitch (alineado a la r\'ubrica)}\\[2pt]
{\large Guion de cabina · 4 minutos · Educat\'on Challenge 2026}
\end{center}
\vspace{2pt}
{\color{yariblue}\hrule height 1pt}
\vspace{5pt}

\bloque{[0:00--0:25] \ HOOK + DATO}{}{%
``Buenas tardes. Una pregunta: ¿qu\'e pasa cuando 30 estudiantes quieren cambiar algo en su colegio\dots pero nadie los modera, se basan en un TikTok falso y terminan gritando sin llegar a nada? En el Per\'u, \textbf{solo 1 de cada 3 estudiantes de secundaria entiende qu\'e es la democracia}; y aunque la mitad vota por sus representantes, \textbf{solo 2 de cada 10 participan en decisiones reales}.''}

\bloque{[0:25--0:50] \ MARIANA + PROBLEMA REAL}{Impacto social}{%
``Entrevistamos a Mariana, 16 a\~nos. Nos dijo lo que lo resume todo: \textit{`S\'i tenemos el espacio\dots lo m\'as dif\'icil es ponernos de acuerdo, porque cada uno tiene su opini\'on y chocan.'} El problema no es que no quieran participar. Es que ante un \textbf{desaf\'io real de su comunidad}, el debate se vuelve caos ---o se quedan callados por miedo---, y \textbf{aprenden que participar no sirve}.''}

\bloque{[0:50--1:10] \ EJE + VISI\'ON}{Pertinencia tem\'atica}{%
``Nuestra propuesta se enmarca en el eje de \textbf{Innovaci\'on para la participaci\'on ciudadana}. Porque la C\'ivica \textbf{no es debatir para ganar: es deliberar para comprender un problema y construir consenso}. Con sustento pedag\'ogico: la ciudadan\'ia se aprende ejerci\'endola, no memoriz\'andola.''}

\bloque{[1:10--1:50] \ LA SOLUCI\'ON: 3 CAMBIOS}{Originalidad e innovaci\'on}{%
``Presentamos \textbf{YariNET}, y es original en todos sus componentes, \textbf{adaptada al aula peruana}: cambiamos \textbf{el espacio} ---el aula en c\'irculo, sin jerarqu\'ias---; \textbf{el dispositivo} ---el celular que ya tienen, con un soporte que imprimimos en 3D---; y ponemos \textbf{una IA} que ---junto al profesor--- \textbf{no impone: educa}. Somos la \'unica propuesta que integra las tres cosas, en espa\~nol y sobre el curso que ya existe.''}

\bloque{[1:50--2:30] \ C\'OMO FUNCIONA (mostrar soporte)}{Impacto social}{%
``Miren. \textit{[muestras el soporte 3D con el celular]} El estudiante \textbf{identifica} un problema real de su comunidad, \textbf{delibera} con sus compa\~neros ---la IA modera con imparcialidad, verifica las fuentes y marca las fake news, y sintetiza el consenso---, y \textbf{ejecuta} una propuesta que llega al director. \textbf{Identificar, deliberar, ejecutar}: eso es un agente de cambio de verdad.''}

\bloque{[2:30--3:05] \ VIABILIDAD Y ESCALABILIDAD}{Viabilidad (clave)}{%
``¿Es viable? \textbf{No es una idea en papel: ya lo tenemos funcionando.} Usa lo que el colegio \textbf{ya tiene} ---el celular y el curso de C\'ivica---; cuesta \textbf{centavos por estudiante}, del orden de 5 d\'olares al mes por aula. Y es \textbf{replicable en cualquier colegio del Per\'u, y escalable a universidades}.''}

\bloque{[3:05--3:30] \ EVIDENCIA + DIFERENCIADOR}{Originalidad / Viabilidad}{%
``¿Y funciona el m\'etodo? \textbf{Google DeepMind lo public\'o en Science en 2024}: una IA ayud\'o a grupos a acordar mejor que un mediador humano. A diferencia de otros que solo ofrecen \textit{`un espacio para dialogar'}, \textbf{nosotros damos la moderaci\'on que hace que el acuerdo realmente ocurra}.''}

\bloque{[3:30--4:00] \ CIERRE / SLOGAN}{Impacto social}{%
``YariNET forma \textbf{agentes de cambio de verdad}: que identifican, deliberan y ejecutan; que \textbf{no le temen al debate} y aun as\'i construyen acuerdos. Porque la ciudadan\'ia no se memoriza ---\textbf{se ejerce}. {\color{yarigreen}\textbf{YariNET: constructores de acuerdos.}} Gracias.''}

\vspace{4pt}
{\color{yariblue}\hrule height 1pt}
\vspace{4pt}

\noindent{\color{yariblue}\bfseries Mapeo a la r\'ubrica:} \
\textbf{Originalidad 25\%} (bloques 4 y 6) · \textbf{Viabilidad 25\%} (bloque 6) · \textbf{Pertinencia 20\%} (bloque 3) · \textbf{Impacto social 20\%} (bloques 2, 5 y 8) · \textbf{Comunicaci\'on 10\%} (estructura + demo visual).

\vspace{2pt}
\noindent{\color{yariblue}\bfseries Palabras que el jurado quiere o\'ir:} \ ``viable y ejecutable'', ``replicable en colegios y universidades'', ``adaptada al entorno escolar local'', ``eje de Innovaci\'on para la participaci\'on ciudadana'', ``agente de cambio que identifica, delibera y ejecuta''.

\vspace{2pt}
\noindent{\color{yariblue}\bfseries Tips (d\'ia D):} \ Prioriza \textbf{Viabilidad + el soporte 3D + el cierre}. \textbf{No presentes al equipo}. Ten la \textbf{captura de la app} en la diapositiva de atr\'as (refuerza ``ya funciona'').

\end{document}
